\documentclass[11pt]{article}
\usepackage{amsmath,amssymb}
\usepackage[margin=1in]{geometry}

\title{The Cayley Transform Plays Music}
\author{}
\date{}

\begin{document}
\maketitle

\noindent\textbf{Hook:} The octave, the fifth, and the fourth walk into a bar.
They are $Q(1/3)$, $Q(1/5)$, and $Q(1/7)$.

\medskip

The \emph{Cayley transform} $Q(x) = \frac{1+x}{1-x}$ is a M\"obius transformation
that converts between bounded and unbounded quantities --- between velocities and
energies in relativity, between tanh and exp in analysis, between tournament parameters
and path counts in combinatorics.

Feed it reciprocal odd numbers and something unexpected happens.

\begin{theorem}
For every $n \geq 1$,
\[
Q\!\left(\frac{1}{2n+1}\right) = \frac{n+1}{n}.
\]
\end{theorem}

\begin{proof}
$Q(1/(2n+1)) = \frac{1 + 1/(2n+1)}{1 - 1/(2n+1)} = \frac{2n+2}{2n} = \frac{n+1}{n}$.\qed
\end{proof}

The ratios $(n+1)/n$ are the \emph{superparticular ratios}, the foundation of
just intonation in Western music theory:
\begin{center}
\begin{tabular}{ccl}
$n$ & $Q(1/(2n+1))$ & Musical interval \\
\hline
1 & $2/1$ & octave \\
2 & $3/2$ & perfect fifth \\
3 & $4/3$ & perfect fourth \\
4 & $5/4$ & major third \\
5 & $6/5$ & minor third \\
\end{tabular}
\end{center}

\noindent The inverse is equally clean: $Q^{-1}\!\left(\frac{n+1}{n}\right) = \frac{1}{2n+1}$.
Each musical interval \emph{encodes} an odd number. Consonance decreases as the
odd number increases: the octave ($1/3$) is more consonant than the fifth ($1/5$),
which is more consonant than the fourth ($1/7$), and so on.

The connection to analysis: the Taylor series of the inverse hyperbolic tangent is
\[
\operatorname{arctanh}(x) = x + \frac{x^3}{3} + \frac{x^5}{5} + \frac{x^7}{7} + \cdots
\]
The denominators $1, 3, 5, 7, \ldots$ are the same odd numbers. Through $Q$, each
Taylor coefficient $1/(2n+1)$ maps to the musical interval $(n+1)/n$. The
\textbf{harmonic series of arctanh is the harmonic series of just intonation}.

Meanwhile, $Q(x) = e^{2\operatorname{arctanh}(x)}$, so multiplication of $Q$-values
corresponds to addition of arctanh distances. Two successive musical intervals
\emph{multiply} their frequency ratios, which \emph{adds} their arctanh
representations. An octave plus a fifth is $Q(1/3) \cdot Q(1/5) = 2 \cdot 3/2 = 3$,
which equals $Q(1/2)$ --- the perfect twelfth.

\medskip

\noindent\textbf{Bonus:} $Q$ also cubes the golden ratio ($Q(1/\varphi) = \varphi^3$)
and squares the tribonacci constant ($Q(1/\tau) = \tau^2$). It shifts Fibonacci
ratios by three indices: $Q(F_n/F_{n+1}) = F_{n+2}/F_{n-1}$. And the product
$\prod_{p \leq 19} Q(1/p)$ over the first eight primes equals $21 = 3 \times 7$
--- the forbidden Hamiltonian path count in tournament theory.

\end{document}
